% Source: https://artofproblemsolving.com/community/c3h2267005_probability_question
% Stable identity: aops-topic-2267005-q01
\textbf{(a)} By independence, the probability is
\[P_k=p^k\left(\dfrac{1}{5}\right)^{5-k}=\left(\dfrac{1}{5}\right)^5(5p)^k.\]

\textbf{(b)} As $k$ goes from $0$ to $5$, $P_k$ is multiplied by $5p$ at each step.
\begin{itemize}
\item If $p>\dfrac{1}{5}$, then $5p>1$, so $P_k$ increases with $k$: he should \textbf{reason on all five} questions, with chance $p^5$.
\item If $p<\dfrac{1}{5}$, then $5p<1$, so $P_k$ decreases: he should \textbf{guess on all five}, with chance $\left(\dfrac{1}{5}\right)^5=\dfrac{1}{3125}$.
\item If $p=\dfrac{1}{5}$, all choices of $k$ give the same chance $\dfrac{1}{3125}$.
\end{itemize}
Since $P_k$ is always monotonic in $k$, the largest value is at $k=0$ or $k=5$: \textbf{a combination is never strictly better} than the better of the two pure strategies.
